% titel: Blatt 7, A5: Partialdrücke in einem Gasgemisch (He, CH4, C2H6)
% skript: Mortimer Kap. 11 (ideale Gase, Daltonsches Gesetz)

% Teil I: Werkzeuge, kein L\"osungsweg.

In einem idealen Gasgemisch trägt jede Komponente unabh\"angig von den
anderen zum Gesamtdruck bei.

\begin{werkzeugkasten}[Daltonsches Partialdruckgesetz]
F\"ur ein ideales Gasgemisch gilt:
\[
  p_i = n_i\,\frac{RT}{V} , \qquad
  p_{\text{ges}} = \sum_i p_i .
\]
Mit Stoffmenge $n_i = m_i/M_i$ und Gasgleichung $pV = nRT$.
Jede Komponente ``f\"ullt'' das gesamte Volumen $V$ aus\footnote{F\"ur
ideale Gase ist das Mischvolumen die Summe der Einzelvolumina beim selben
$p,T$; der Partialdruck ist proportional zum Stoffmengenanteil
$x_i = n_i/n_{\text{ges}}$, d.~h.\ $p_i = x_i\,p_{\text{ges}}$.}.
\textit{Skript:} Mortimer Kap. 11, ideales Gas und Daltonsches Gesetz.
\end{werkzeugkasten}

\begin{achtung}[Darauf achten]
\begin{itemize}
  \item Stoffmenge aus Masse und Molmasse: Methan $\mathrm{CH_4}$ hat
        $M=16.04$, Ethan $\mathrm{C_2H_6}$ $M=30.07$, Helium $M=4.00$.
  \item Temperaturen in Kelvin (hier 300 K) einsetzen.
  \item Gesamtdruck $=$ Summe der Partialdr\"ucke; Mohrwerte runden.
\end{itemize}
\end{achtung}
